\begin{thebibliography}{99}

\bibitem{lennie2003cost} Lennie, P. The cost of cortical computation. \emph{Current Biology}, \textbf{13}, 493--497 (2003).

\bibitem{31} Niven, J.E. \& Laughlin, S.B. Energy limitation as a selective pressure on the evolution of sensory systems. \emph{Journal of Experimental Biology}, \textbf{211}, 1792--1804 (2008).

\bibitem{32} Niven, J.E. Neuronal energy consumption: Biophysics, efficiency and evolution. \emph{Current Opinion in Neurobiology}, \textbf{41}, 129--135 (2016).

\bibitem{1} Rolls, E.T. \& Tovee, M.J. Sparseness of the neuronal representation of stimuli in the primate temporal visual cortex. \emph{Journal of Neurophysiology}, \textbf{73}, 713--726 (1995).

\bibitem{2} Treves, A. \& Rolls, E. What determines the capacity of autoassociative memories in the brain?. \emph{Network: Computation in Neural Systems}, \textbf{2}, 371--397 (2009).

\bibitem{3} Haider, B., Krause, M.R., Duque, A., Yu, Y., Touryan, J., Mazer, J.A. \& McCormick, D.A. Synaptic and network mechanisms of sparse and reliable visual cortical activity during nonclassical receptive field stimulation. \emph{Neuron}, \textbf{65}, 107--121 (2010).

\bibitem{20} Vinje, W.E. \& Gallant, J.L. Sparse coding and decorrelation in primary visual cortex during natural vision. \emph{Science}, \textbf{287}, 1273--1276 (2000).

\bibitem{30} Barlow, H.B. Possible principles underlying the transformations of sensory messages. MIT Press, 1961.

\bibitem{13} Olshausen, B.A. \& Field, D.J. Sparse coding of sensory inputs. \emph{Current Opinion in Neurobiology}, \textbf{14}, 481--487 (2004).

\bibitem{field_relations_1987} Field, D.J. Relations between the statistics of natural images and the response properties of cortical cells. \emph{Journal of the Optical Society of America A}, \textbf{4}, 2379 (1987).

\bibitem{hubel_receptive_1962} Hubel, D.H. \& Wiesel, T.N. Receptive fields, binocular interaction and functional architecture in the cat's visual cortex. \emph{Journal of Physiology}, \textbf{160}, 106--154 (1962).

\bibitem{15} Franco, L., Rolls, E.T., Aggelopoulos, N.C. \& Jerez, J.M. Neuronal selectivity, population sparseness, and ergodicity in the inferior temporal visual cortex. \emph{Biological Cybernetics}, \textbf{96}, 547--560 (2007).

\bibitem{22} Yu, Y., McTavish, T.S., Hines, M.L., Shepherd, G.M., Valenti, C. \& Migliore, M. Sparse distributed representation of odors in a large-scale olfactory bulb circuit. \emph{PLoS Computational Biology}, \textbf{9}, e1003014 (2013).

\bibitem{25} Lin, A.C., Bygrave, A.M., de Calignon, A., Lee, T. \& Miesenbock, G. Sparse, decorrelated odor coding in the mushroom body enhances learned odor discrimination. \emph{Nature Neuroscience}, \textbf{17}, 559--568 (2014).

\bibitem{6} Ecker, A.S., Berens, P., Tolias, A.S. \& Bethge, M. The effect of noise correlations in populations of diversely tuned neurons. \emph{Journal of Neuroscience}, \textbf{31}, 14272--14283 (2011).

\bibitem{4} Softky, W.R. \& Koch, C. The highly irregular firing of cortical cells is inconsistent with temporal integration of random EPSPs. \emph{Journal of Neuroscience}, \textbf{13}, 334--350 (1993).

\bibitem{26} Brown, D., Feng, J. \& Feerick, S. Variability of firing of Hodgkin-Huxley and FitzHugh-Nagumo neurons with stochastic synaptic input. \emph{Physical Review Letters}, \textbf{82}, 4731--4734 (1999).

\bibitem{27} Salinas, E. \& Sejnowski, T.J. Impact of correlated synaptic input on output firing rate and variability in simple neuronal models. \emph{Journal of Neuroscience}, \textbf{20}, 6193--6209 (2000).

\bibitem{28} Feng, J. \& Brown, D. Coefficient of variation of interspike intervals greater than 0.5: How and when?. \emph{Biological Cybernetics}, \textbf{80}, 291--297 (1999).

\bibitem{39} Stiefel, K.M., Englitz, B. \& Sejnowski, T.J. Origin of intrinsic irregular firing in cortical interneurons. \emph{Proceedings of the National Academy of Sciences}, \textbf{110}, 7886--7891 (2013).

\bibitem{7} Marder, E. Variability, compensation, and modulation in neurons and circuits. \emph{Proceedings of the National Academy of Sciences}, \textbf{108}, 15542--15548 (2011).

\bibitem{8} Faisal, A.A., Selen, L.P. \& Wolpert, D.M. Noise in the nervous system. \emph{Nature Reviews Neuroscience}, \textbf{9}, 292--303 (2008).

\bibitem{9} Festa, D., Aschner, A., Davila, A., Kohn, A. \& Coen-Cagli, R. Neuronal variability reflects probabilistic inference tuned to natural image statistics. \emph{Nature Communications}, \textbf{12}, 3635 (2021).

\bibitem{5} Lengler, J. \& Steger, A. Note on the coefficient of variations of neuronal spike trains. \emph{Biological Cybernetics}, \textbf{111}, 229--235 (2017).

\bibitem{47} Manwani, A., Steinmetz, P.N. \& Koch, C. The impact of spike timing variability on the signal-encoding performance of neural spiking models. \emph{Neural Computation}, \textbf{14}, 347--367 (2002).

\bibitem{10} Abdi, H. Coefficient of variation. \emph{Encyclopedia of Statistical Sciences}, \textbf{94}, 94 (2011).

\bibitem{11} Miyawaki, H., Watson, B.O. \& Diba, K. Neuronal firing rates diverge during REM and homogenize during non-REM. \emph{Scientific Reports}, \textbf{9}, 689 (2019).

\bibitem{12} Rudolph, M. \& Destexhe, A. The discharge variability of neocortical neurons during high-conductance states. \emph{Neuroscience}, \textbf{119}, 855--873 (2003).

\bibitem{Buracas1998} Bura{\v c}as, G.T., Zador, A.M., DeWeese, M.R. \& Albright, T.D. Efficient discrimination of temporal patterns by motion-sensitive neurons in primate visual cortex. \emph{Neuron}, \textbf{20}, 959--969 (1998).

\bibitem{Renart2010} Renart, A., de la Rocha, J., Bartho, P., Hollender, L., Parga, N., Reyes, A. \& Harris, K.D. The asynchronous state in cortical circuits. \emph{Science}, \textbf{327}, 587--590 (2010).

\bibitem{LitwinKumar2012} Litwin-Kumar, A. \& Doiron, B. Slow dynamics and high variability in balanced cortical networks with clustered connections. \emph{Nature Neuroscience}, \textbf{15}, 1498--1505 (2012).

\bibitem{41} Sukenik, N., Vinogradov, O., Weinreb, E., Segal, M., Levina, A. \& Moses, E. Neuronal circuits overcome imbalance in excitation and inhibition by adjusting connection numbers. \emph{Proceedings of the National Academy of Sciences}, \textbf{118}, e2018459118 (2021).

\bibitem{Rieke1997} Rieke, F., Warland, D., de Ruyter van Steveninck, R. \& Bialek, W. Spikes: Exploring the Neural Code. MIT Press, 1997.

\bibitem{Sterling2015} Sterling, P. \& Laughlin, S. Principles of Neural Design. MIT Press, 2015.

\bibitem{Laughlin1981} Laughlin, S.B. A Simple Coding Procedure Enhances a Neuron's Information Capacity. \emph{Zeitschrift für Naturforschung C}, \textbf{36}, 910--912 (1981).

\bibitem{Bell1995Infomax} Bell, A.J. \& Sejnowski, T.J. An information-maximization approach to blind separation and blind deconvolution. \emph{Neural Computation}, \textbf{7}, 1129--1159 (1995).

\bibitem{Yu2005VisionResearch} Yu, Y., Potetz, B. \& Lee, T.S. The role of spiking nonlinearity in contrast gain control and information transmission. \emph{Vision Research}, \textbf{45}, 583--592 (2005).

\bibitem{45} Tkacik, G., Prentice, J.S., Balasubramanian, V. \& Schneidman, E. Optimal population coding by noisy spiking neurons. \emph{Proceedings of the National Academy of Sciences}, \textbf{107}, 14419--14424 (2010).

\bibitem{Chalk2018} Chalk, M., Marre, O. \& Tka{\v{c}}ik, G. Toward a unified theory of efficient, predictive, and sparse coding. \emph{Proceedings of the National Academy of Sciences}, \textbf{115}, 186--191 (2018).

\bibitem{Raichle2001} Raichle, M.E., MacLeod, A.M., Snyder, A.Z., Powers, W.J., Gusnard, D.A. \& Shulman, G.L. A default mode of brain function. \emph{Proceedings of the National Academy of Sciences}, \textbf{98}, 676--682 (2001).

\bibitem{Tomasi2013} Tomasi, D., Wang, G. \& Volkow, N.D. Energetic cost of brain functional connectivity. \emph{Proceedings of the National Academy of Sciences}, \textbf{110}, 13642--13647 (2013).

\bibitem{23} Semedo, J.D., Zandvakili, A., Machens, C.K., Yu, B.M. \& Kohn, A. Cortical areas interact through a communication subspace. \emph{Neuron}, \textbf{102}, 249--259 (2019).

\bibitem{Jia2024variability} Akella, S., Ledochowitsch, P., Siegle, J.H., Belski, H., Denman, D.J., Buice, M.A., Durand, S., Koch, C., Olsen, S.R. \& Jia, X. Deciphering neuronal variability across states reveals dynamic sensory encoding. \emph{bioRxiv} (2024).

\bibitem{strong1998entropy} Strong, S.P., Koberle, R., de Ruyter van Steveninck, R.R. \& Bialek, W. Entropy and information in neural spike trains. \emph{Physical Review Letters}, \textbf{80}, 197--200 (1998).

\bibitem{16} Vogels, T.P., Sprekeler, H., Zenke, F., Clopath, C. \& Gerstner, W. Inhibitory plasticity balances excitation and inhibition in sensory pathways and memory networks. \emph{Science}, \textbf{334}, 1569--1573 (2011).

\bibitem{sp_stdp} Kim, H., Tang, H., Choi, W. \& Park, J. An energy-quality scalable STDP-based sparse coding processor with on-chip learning capability. \emph{IEEE Transactions on Biomedical Circuits and Systems}, \textbf{14}, 125--137 (2020).

\bibitem{brendel_learning_2020} Brendel, W., Bourdoukan, R., Vertechi, P., Machens, C.K. \& Den{\`e}ve, S. Learning to represent signals spike by spike. \emph{PLoS Computational Biology}, \textbf{16}, e1007692 (2020).

\bibitem{Sp2019} Woods, W. \& Teuscher, C. Fast and accurate sparse coding of visual stimuli with a simple, ultralow-energy spiking architecture. \emph{IEEE Transactions on Neural Networks and Learning Systems}, \textbf{30}, 2173--2187 (2019).

\bibitem{Yuguo2018Evaluating} Yu, Y., Herman, P., Rothman, D.L., Agarwal, D. \& Hyder, F. Evaluating the gray and white matter energy budgets of human brain function. \emph{Journal of Cerebral Blood Flow and Metabolism}, \textbf{38}, 1339--1353 (2018).

\bibitem{37} Olshausen, B.A. \& Field, D.J. Emergence of simple-cell receptive field properties by learning a sparse code for natural images. \emph{Nature}, \textbf{381}, 607--609 (1996).

\bibitem{38} Olshausen, B.A. \& Field, D.J. Sparse coding with an overcomplete basis set: A strategy employed by V1?. \emph{Vision Research}, \textbf{37}, 3311--3325 (1997).

\bibitem{97bialek} de Ruyter van Steveninck, R.R., Lewen, G.D., Strong, S.P., Koberle, R. \& Bialek, W. Reproducibility and variability in neural spike trains. \emph{Science}, \textbf{275}, 1805--1808 (1997).

\bibitem{wang2004image} Wang, Z., Bovik, A.C., Sheikh, H.R. \& Simoncelli, E.P. Image quality assessment: from error visibility to structural similarity. \emph{IEEE Transactions on Image Processing}, \textbf{13}, 600--612 (2004).

\bibitem{17} Hartmann, K., Bruehl, C., Golovko, T. \& Draguhn, A. Fast homeostatic plasticity of inhibition via activity-dependent vesicular filling. \emph{PLoS One}, \textbf{3}, e2979 (2008).

\bibitem{18} Woodin, M.A., Ganguly, K. \& Poo, M. Coincident pre- and postsynaptic activity modifies GABAergic synapses by postsynaptic changes in Cl- transporter activity. \emph{Neuron}, \textbf{39}, 807--820 (2003).

\bibitem{laughlin_energy_2001} Laughlin, S. Energy as a constraint on the coding and processing of sensory information. \emph{Current Opinion in Neurobiology}, \textbf{11}, 475--480 (2001).

\bibitem{Harris2012} Harris, J.J., Jolivet, R. \& Attwell, D. Synaptic energy use and supply. \emph{Neuron}, \textbf{75}, 762--777 (2012).

\bibitem{Shoval2012} Shoval, O., Sheftel, H., Shinar, G., Hart, Y., Ramote, O., Mayo, A., Dekel, E. \& Alon, U. Evolutionary trade-offs, Pareto optimality, and the geometry of phenotype space. \emph{Science}, \textbf{336}, 1157--1160 (2012).

\bibitem{Sommer2026} Sommer, P., Schumacher, J. \& Gjorgjieva, J. Neuronal degeneracy in information–energy trade-offs. \emph{bioRxiv} (2026).

\bibitem{friston_free-energy_2010} Friston, K. The free-energy principle: A unified brain theory?. \emph{Nature Reviews Neuroscience}, \textbf{11}, 127--138 (2010).

\bibitem{43} Shew, W.L., Yang, H., Yu, S., Roy, R. \& Plenz, D. Information capacity and transmission are maximized in balanced cortical networks with neuronal avalanches. \emph{Journal of Neuroscience}, \textbf{31}, 55--63 (2011).

\bibitem{Liang2025prl} Yang, Z., Liang, J. \& Zhou, C. Critical avalanches in excitation-inhibition balanced networks reconcile response reliability with sensitivity for optimal neural representation. \emph{Physical Review Letters}, \textbf{134}, 028401 (2025).

\end{thebibliography}